\section{Circle regularization and transfer from the tube}\label{sec:tube}

Let \(T_\delta=\{z:\operatorname{dist}(z,\Gamma)\le\delta\}\). In
\eqref{eq:contour-admissible} replace only the support condition on \(\nu\)
by \(\operatorname{supp}\nu\subset T_\delta\). For \(0<\delta\le1/64\) we
have the exact estimate
\begin{equation}\label{eq:tube-transfer}
 \inf_{T_\delta}\mathcal E\ge\alpha_\Gamma-300\delta.
\end{equation}

The proof covers the unrestricted contour self-energy. Let \(D\) be the
interior of the ellipse. For a point \(z\) inside \(D\), let \(\omega_z\) be
its interior harmonic measure on \(\Gamma\). Outside \(D\), use exterior
harmonic measure with total mass one. At a boundary point use
\(\delta_z\). Set
\[
 c=(A+B)/2,\quad d=(A-B)/2,\quad
 F(w)=\xi+cw+d/w\quad(|w|\ge1).
\]
The factorization \(F(w_1)-F(w_2)=(w_1-w_2)(c-d/(w_1w_2))\) shows that
\(F\) is one-to-one outside the unit disk. Let \(W\) be its inverse. Set
\(e(z)=0\) inside the ellipse and \(e(z)=\log|W(z)|\) outside. The
Green-function identity gives
\begin{equation}\label{eq:tube-green}
 U^{\omega_z}(x)\le K_0(z,x)+e(z),\qquad
 I(\omega_z,\omega_w)\le K_0(z,w)+e(z)+e(w).
\end{equation}
For two interior points the first deficit is the interior Green function.
For points on opposite sides it is zero. For two exterior points it is the
nonnegative Green function of the exterior of the disk,
\[
 \log\left|\frac{W(z)\overline{W(x)}-1}{W(z)-W(x)}\right|.
\]
Its sign follows from
\(|a\overline b-1|^2-|a-b|^2=(|a|^2-1)(|b|^2-1)\ge0\). At boundary points
the identity in \eqref{eq:tube-green} is an equality. Integrating it
against \(\omega_w\) proves the two-sweep inequality. These identities also
follow by Dirichlet uniqueness after subtracting the Green pole, including
the growth term at infinity. Integrable logarithms and limiting arguments
cover boundary points and arbitrary finite-energy measures.

If \(z\in T_\delta\) is exterior, the factorization gives
\(\operatorname{dist}(z,\Gamma)\ge B(|W(z)|-1)\). Hence
\begin{equation}\label{eq:tube-e}
 0\le e(z)\le\delta/B.
\end{equation}

Let \(q(z)=1-(\Re z-\xi)^2/A^2-(\Im z)^2/B^2\). On \(D\),
\(|\nabla q|\le2/B\). For \(x\in[0,2]\), \(q(x)>1/5\). A point in the inner
\(\delta\)-collar has \(q(z)\le2\delta/B<1/8\). The level set \(q=1/8\) has
distance more than \(1/32\) from \([0,2]\). By domain monotonicity the
interior Green function is at most \(\log160<6\) there. The superharmonic
barrier \(48q\) dominates it on both boundaries of the collar. By the
maximum principle,
\begin{equation}\label{eq:tube-gD}
 g_D(z,x)\le96\delta/B
       \quad(z\in D\cap T_\delta,\ x\in[0,2]).
\end{equation}

Put \(f=A_0-\phi\). Let \(z\) be interior. Sweeping the positive-source
logarithms in \(A_0\) increases their sum by at most \(3\cdot96\delta/B\).
The linear part of \(\phi\) keeps its mean. By \eqref{eq:tube-green}, the
logarithmic-potential part of \(\phi\) has the favorable sign in \(-\phi\).
Thus \(\int f\,d\omega_z-f(z)\le288\delta/B\). Now let \(z=F(w)\) be
exterior. Its mean coordinate is \(\xi+c/\overline w+d/w\), at distance at
most \((A+B)(|w|-1)\) from \(z\). The change in \(A_0\) is exactly
\(-3e(z)\). The remaining terms give the smaller bound
\([1/2+(A+B)\log2]\delta/B<3\delta/B\). Hence, for the swept measure
\(\nu^\#=\int\omega_z\,d\nu(z)\) of mass \(s\),
\[
 \int f\,d\nu^\#\le\int f\,d\nu+288s\delta/B.
\]

Put \(b=\int e\,d\nu\). The changed interaction energies in
\eqref{eq:contour-energy} increase by at most
\(2sb+2(1-s)b+b=3b\le3s\delta/B\). Hence
\[
 \mathcal E(\kappa+\nu^\#,\mu)
  \le\mathcal E(\kappa+\nu,\mu)+(291/B)s\delta
  <\mathcal E(\kappa+\nu,\mu)+300\delta.
\]
This proves \eqref{eq:tube-transfer}, uniformly for every frequency measure
on the full half-line and every allowed pinned/contour mass split.

Let \(z_1,\ldots,z_d\in\Gamma\). Let \(\sigma_{z,\delta}\) be probability
arclength on the circle of radius \(\delta\) about \(z\). Set
\(\nu_\delta=m^{-1}\sum_{a=1}^d\sigma_{z_a,\delta}\). The circle mean
identity \(\int\log|x-z-\delta e^{iu}|\,du/(2\pi)=\log\max(|x-z|,\delta)\)
implies
\begin{equation}\label{eq:tube-circle}
 \frac1{m^2}\sum_{a\ne b}K_0(z_a,z_b)
 \ge I(\nu_\delta)-\frac d{m^2}\log(1/\delta).
\end{equation}
The mixed potentials from \([0,2]\) and \((-\infty,-1/2]\) are harmonic on
these disks. The function \(\phi\) is subharmonic, also at the interior of
its zero segment. So the circular mean of the contour field \(f\) is at
most its central value. Take \(\delta=m^{-2}\) and \(m\ge8\). Then
\eqref{eq:tube-circle} and \eqref{eq:tube-transfer} together cost at most
\(2d\log m+300\) in the unnormalized energy. Coincident particles give a
zero Vandermonde, so no separation hypothesis is needed. A zero of \(F_g\)
at finite \(m\) only decreases the absolute integrand. No two-sided
logarithmic estimate at finite \(m\) is used.
